% ============================================================
%  CHAPTER 3 — SOFTWARE REQUIREMENTS SPECIFICATION (SRS)
% ============================================================
\chapter{Software Requirements Specification}


\section{Overall Description}

\subsection{Product Perspective}

MindTrace is a standalone web-based application that integrates a deep
learning inference backend with a browser-based frontend. It does not depend
on third-party emotion analysis APIs; all inference is performed by the
custom-trained ResNet-18 model hosted on the cloud. The system interacts with
a MongoDB database for data persistence and with the client's webcam for
real-time frame capture.

\subsection{Product Functions}

\begin{itemize}
    \item Real-time facial emotion classification into seven categories.
    \item Head pose estimation to detect attention direction.
    \item Blink-rate measurement for fatigue detection.
    \item Composite engagement score computation (0--100 scale).
    \item Session-wise logging of emotion and engagement data.
    \item Interactive analytics dashboard with charts and timelines.
    \item Role-based user management (student and administrator roles).
    \item Secure user authentication with JWT and bcrypt.
\end{itemize}

\subsection{User Classes and Characteristics}

\begin{itemize}
    \item \textbf{Students / Employees}: Use the system during sessions to
          have their engagement monitored. They can view their own session
          history and emotion analytics through the dashboard.
    \item \textbf{Administrators / Instructors}: Have elevated access to
          view aggregated engagement reports for all users under their
          purview and manage user accounts.
\end{itemize}

\subsection{Operating Environment}

The system operates in the following environment:
\begin{itemize}
    \item \textbf{Backend}: Python 3.10+, FastAPI, deployed on Hugging Face
          Spaces (Linux container).
    \item \textbf{Database}: MongoDB Atlas (cloud-hosted).
    \item \textbf{Frontend}: React + Vite, deployed on Vercel.
    \item \textbf{Client}: Any modern web browser (Chrome, Firefox, Edge)
          with webcam access.
\end{itemize}

\section{Functional Requirements}

\subsection{User Authentication}
\begin{itemize}
    \item The system shall allow users to register with a username, email,
          password, and role.
    \item Passwords shall be hashed using bcrypt before storage.
    \item On login, the system shall issue a signed JWT token valid for a
          configurable duration.
    \item Protected endpoints shall reject requests with invalid or expired
          tokens with an HTTP 401 response.
\end{itemize}

\subsection{Emotion Detection}
\begin{itemize}
    \item The system shall accept a Base64-encoded webcam frame via a REST
          endpoint.
    \item It shall detect faces using MediaPipe Face Mesh and extract the
          face region with padding.
    \item The extracted face region shall be classified by the ResNet-18
          model into one of seven emotional states.
    \item The API response shall include the detected emotion label and its
          confidence score.
\end{itemize}

\subsection{Engagement Scoring}
\begin{itemize}
    \item The system shall estimate head pose (pitch, yaw, roll) using a
          3D-2D point correspondence (PnP) solver.
    \item It shall calculate blink rate from eye-aspect-ratio landmarks.
    \item A composite engagement score shall be computed by combining
          emotion weight, head pose deviation, and blink frequency.
    \item The score shall be updated and returned with each API call.
\end{itemize}

\subsection{Session Management}
\begin{itemize}
    \item The system shall persist each session's start time, end time,
          per-frame emotion logs, and computed engagement scores in MongoDB.
    \item Users shall be able to retrieve a history of their past sessions.
    \item Administrators shall be able to view session data for all users.
\end{itemize}

\subsection{Analytics Dashboard}
\begin{itemize}
    \item The frontend shall display a live emotion label and engagement
          score overlay on the webcam feed.
    \item It shall render a timeline chart of engagement scores for the
          current session.
    \item Emotion distribution shall be shown as a pie/bar chart.
    \item Session history shall be accessible in a paginated table view.
\end{itemize}

\section{Non-Functional Requirements}

\subsection{Performance}
\begin{itemize}
    \item The system shall process a single webcam frame (inference +
          engagement scoring) within \textbf{300 ms} under normal load.
    \item The frontend shall update the emotion overlay at a minimum of
          \textbf{5 frames per second}.
\end{itemize}

\subsection{Scalability}
\begin{itemize}
    \item The stateless backend design shall allow horizontal scaling by
          deploying multiple FastAPI instances behind a load balancer.
    \item MongoDB Atlas shall support automatic scaling of storage.
\end{itemize}

\subsection{Security}
\begin{itemize}
    \item All API communication shall occur over HTTPS.
    \item JWT tokens shall use HS256 signing and include an expiry claim.
    \item User passwords shall never be stored in plain text.
\end{itemize}

\subsection{Usability}
\begin{itemize}
    \item The user interface shall be operable without installation of any
          additional software or browser extensions.
    \item The dashboard shall be responsive and usable on screen widths
          from 768 px (tablet) upwards.
\end{itemize}

\subsection{Reliability}
\begin{itemize}
    \item The system shall gracefully handle frames in which no face is
          detected, returning a neutral engagement state without crashing.
    \item The frontend shall display a user-friendly error message if the
          backend is unreachable.
\end{itemize}

\section{System Requirements}

\subsection{Hardware Requirements}

\begin{table}[H]
\centering
\caption{Hardware Requirements}
\begin{tabular}{ll}
\toprule
\textbf{Component} & \textbf{Minimum Specification} \\
\midrule
Processor       & Intel Core i5 / AMD Ryzen 5 (or equivalent) \\
RAM             & 8 GB \\
Webcam          & 720p USB or built-in camera \\
Network         & Broadband internet connection (5 Mbps+) \\
Storage (client)& 500 MB free disk space (browser cache) \\
\bottomrule
\end{tabular}
\end{table}

\subsection{Software Requirements}

\begin{table}[H]
\centering
\caption{Software Requirements}
\begin{tabular}{ll}
\toprule
\textbf{Component}      & \textbf{Version / Details} \\
\midrule
Python                  & 3.10 or higher \\
PyTorch                 & 2.x (with torchvision) \\
FastAPI                 & 0.110+ \\
MediaPipe               & 0.10+ \\
OpenCV                  & 4.x \\
MongoDB                 & Atlas cloud (or local 6.x) \\
React                   & 18.x \\
Node.js                 & 18 LTS \\
Web Browser             & Chrome 110+ / Firefox 110+ / Edge 110+ \\
\bottomrule
\end{tabular}
\end{table}

\section{Constraints and Assumptions}

\begin{itemize}
    \item The system assumes adequate and consistent lighting for reliable
          face detection. Poor lighting may reduce model accuracy.
    \item The ResNet-18 model is trained exclusively on the RAF-DB dataset;
          performance may vary for demographics underrepresented in that
          dataset.
    \item Browser webcam access requires the user to grant camera
          permissions; denial will prevent system operation.
    \item Latency is subject to network conditions between the client and
          the cloud-hosted backend.
\end{itemize}
